\scanpage{36}

% L1–3
ou encore (pour la 2\textsuperscript{\unsure{ième|ème}} relation)
\[
 \begin{aligned}
 &\lambda(x+u+v)-\lambda(x+u)-\lambda(x+v)+\lambda(x)\\
 &\qquad=\lambda(y+u+v)-\lambda(y+u)-\lambda(y+v)+\lambda(y
 \end{aligned}
\]

% L4–6
\[
 (\Delta_{uv}\lambda)(x)=\mathrm{c}^{\mathrm{te}}\quad(\text{pour }u,v\text{ fixés})
\]
i.e. \deleted{\unsure{Pour poser}}, posant
\[
 \lambda(u+v)-\lambda(u)-\lambda(v)=\mu(u,v)
\]

% L7–8
On doit avoir
\[
 \boxed{\lambda(x+u+v)-\lambda(\unsure{x+u})-\lambda(\unsure{x+v})+\lambda(x)=\mu(u,v)}
\]

% L9–10
i.e. $\mu$ est \uline{bi-additive}.

\deleted{$\lambda(x+u+v)=\lambda(\unsure{u})+\lambda(v)$}

% L11–14
Pour $\mu$ \inserted{(\unsure{Z|2}-bilinéaire symétrique)} \unsure{fixée|fixé}, les $\lambda$ satisfaisant la relation précédente \inserted{et telles que $\lambda(0)=0$} forment un torseur sous l'ens des $f$ telles que

% L14–16
\[
 \left\{\begin{aligned}
 &f(0)=0\\
 &f(\deleted{\unsure{u}}x+u+v)-f(x+u)-f(x+v)\deleted{\unsure{+f(x)}}=0
 \end{aligned}\right.
\]

% L16–19
\inserted{telles que}

\deleted{i.e. $f(x)+f(y)$ ne dépend que de $x+y$, i.e. \unsure{$\exists g:V\to Z$} telle que \unsure{$f(x)+f(y)=f(x+y)$|$f(x+y)+f(0)=f(x$}}

\deleted{\unsure{pour le}} i.e. $f$ est un \deleted{\unsure{f|affine}} \uline{hom}.

% L20–26
\deleted{\unsure{Dans le cas que}} Lorsque $V$ est un module libre sur un $\mathbb Z/n\mathbb Z$, alors \uline{tte} \uline{\unsure{appl.|produit}} bil. sym. de $V$ dans un $Z$ est associée à une forme quadratique (définie mod application \uline{additive}) \unsure{ssi|sss} les $\mu(e_i,e_i)\in2Z$ ($(e_i)$ est une base de $V$ sur $\mathbb Z/n\mathbb Z$) \unsure{en|ou}
% archive sheet number: 35
